\chapter{Hilbert, G\"odel and the Decision Problem}
\label{ch:hilbert}

The theory of computation was born from a question in the foundations of mathematics, and the question whether $\Pc = \NP$ is its direct descendant. This chapter tells that story from Hilbert\index{Hilbert, David}'s program through G\"odel\index{G\"odel, Kurt}'s theorems to the formulation of the Entscheidungsproblem\index{Entscheidungsproblem} that Turing\index{Turing, Alan} answered.

\section{Hilbert's program}

David Hilbert's address to the International Congress of Mathematicians in Paris in 1900 listed, in its published version, twenty three problems, and his tenth asked for a procedure to decide whether a polynomial equation with integer coefficients has integer solutions. The word procedure was not defined; it was understood to mean a finite mechanical rule of the kind a clerk could follow. Hilbert's conviction, stated in the same address and repeated for thirty years, was that every well posed mathematical question has an answer that can be found: there is no ignorabimus.

By 1920 Hilbert had turned this conviction into a program. Mathematics was to be formalised as a system of symbols and rules, the consistency of the system was to be proved by finitary means that no one could doubt, and then the truth of any statement would be a matter of whether it could be derived. The program had three parts: completeness, that every true statement is provable; consistency, that no contradiction is provable; and decidability, that there is a procedure to determine whether a statement is provable. The third is the Entscheidungsproblem, stated precisely by Hilbert and Ackermann in their 1928 textbook for the case of first order logic.

The Entscheidungsproblem is a problem about search. A formula is provable if there exists a finite sequence of formulas, each an axiom or a consequence of earlier ones by a rule, ending in the formula. A proposed proof is easy to check: one verifies each step. The question is whether provability can be decided without searching through all possible proofs, and in this form it is the ancestor of the question whether $\NP = \Pc$. Hilbert expected the answer to be yes, in the strong form that a decision procedure exists; G\"odel's letter of 1956 asks the weaker question whether, for proofs of bounded length, the search can be done in time polynomial in the bound.

\section{G\"odel's theorems}

Kurt G\"odel's paper of 1931 destroyed the first part of the program and the hope of a finitary proof of the second and left the third standing for five more years. The first incompleteness theorem states that any consistent, effectively axiomatized formal system that includes elementary arithmetic contains a statement that is true but not provable in the system. The second states that such a system cannot prove its own consistency. The proof of the first encodes statements and proofs as numbers, G\"odel numbering, and constructs a statement that asserts its own unprovability; if it were provable the system would be inconsistent, so it is unprovable, and hence true.

The technique matters as much as the result. G\"odel showed that the syntax of a formal system, the manipulation of symbols by rules, can be represented within arithmetic, so that statements about proofs are statements about numbers. This is the first instance of what became the central move of computability and complexity theory: the encoding of a computational process as data on which another computational process can act. Turing's universal machine, Cook\index{Cook Levin theorem}'s encoding of a machine's computation as a Boolean formula, and the proof that the queens completion problem is $\NP$ complete all descend from it.

G\"odel's theorems did not settle the Entscheidungsproblem. They showed that provability and truth differ, but the Entscheidungsproblem asked about provability alone, and a decision procedure for provability was not ruled out. It was in fact widely expected that one existed, and the surprise of 1936 was that it did not.

\section{What is a procedure?}

To show that no procedure exists one must say what a procedure is. The three answers given in the mid 1930s were the recursive functions of G\"odel and Herbrand, formalised by Kleene; the lambda calculus of Church\index{Church, Alonzo}; and the machines of Turing. They were shown equivalent within a year or two, and the conjecture that they capture the informal notion of a mechanical procedure is the Church Turing thesis. It is not a theorem, since one side of the equivalence is informal, but it has survived ninety years of attempts to find a counterexample, and every model of computation since proposed, including the random access machine, the cellular automaton, the quantum computer and the neural network with rational weights, computes the same set of functions.

Turing's definition is the one that persuaded. His 1936 paper opens with an analysis of what a human computer does: reads a symbol, changes state of mind, writes a symbol, moves attention to a neighbouring position. The machine that formalises this has a finite set of states, an infinite tape of cells each holding a symbol from a finite alphabet, and a head that reads one cell at a time; a transition table specifies, for each state and symbol, the new state, the symbol to write and the direction to move. Nothing in the definition is arbitrary, and G\"odel himself, who had doubts about the other formalisations, accepted Turing's as decisive.

\section{Church's and Turing's answers}

Church proved in 1936 that there is no lambda definable function that decides whether a formula of first order logic is provable, and Turing proved the same for his machines a few months later, by a different and more fundamental argument. Turing's argument goes through the halting problem: there is no machine that, given the description of a machine $M$ and an input $x$, decides whether $M$ halts on $x$. The proof is a diagonalisation of the kind Cantor used to show the reals uncountable and G\"odel used in the incompleteness theorem. Suppose $H$ decides halting. Build $D$ which, on input a machine description $M$, runs $H$ on $(M, M)$ and halts if and only if $H$ says $M$ does not halt on $M$. Then $D$ on input $D$ halts if and only if it does not. So $H$ does not exist. Since the question whether $M$ halts on $x$ can be written as a formula of first order logic, provable if and only if $M$ halts, a decision procedure for provability would decide halting, and so there is none.

\section{The tenth problem}

Hilbert's tenth problem, on integer solutions of polynomial equations, survived until 1970, when Matiyasevich, building on work of Davis, Putnam and Julia Robinson, showed that every recursively enumerable set is Diophantine: for every set of natural numbers that a Turing machine can enumerate there is a polynomial whose integer solutions, projected on one coordinate, are exactly that set. Since the halting set is recursively enumerable and undecidable, there is no procedure to decide the existence of integer solutions. The theorem is the last of the great negative answers to Hilbert, and it has a complexity theoretic shadow: the existence of a solution to a polynomial system of bounded degree over bounded integers is $\NP$ complete, and the Diophantine encoding of a computation is a cousin of Cook's Boolean encoding.

\section{From computability to complexity}

Turing's theorem draws the line between what can be computed and what cannot. Everything in this book is on the computable side of the line, and the question is not whether but how fast. The step from computability to complexity was taken in the 1960s, when Hartmanis and Stearns defined the time taken by a Turing machine as a function of input length and proved that more time allows more problems to be solved, the time hierarchy theorem of Chapter~\ref{ch:turing}. The proof is a diagonalisation like Turing's, and it is the principal technique known to separate uniform complexity classes unconditionally. That it cannot separate $\Pc$ from $\NP$ is the content of the first of the barriers in Chapter~\ref{ch:barriers}.

The Entscheidungsproblem itself, restricted to proofs of length at most $n$, is in $\NP$ (with the length bound written in unary): the certificate is the proof. It is $\NP$ complete, since any $\NP$ problem can be encoded as the question whether a certain formula has a short proof. G\"odel's letter asked whether it is in $\Pc$. So the question Hilbert asked in 1928, answered negatively by Church and Turing in 1936 for proofs of unbounded length, returns for proofs of bounded length as the question whether $\Pc = \NP$, and it is open.
